%%%%%%%%%%  scan idx 29  |  folio 16  %%%%%%%%%%
\[
f:X\longrightarrow Y
\]
est un
morphisme propre de sch\'emas noeth\'eriens, alors
\begin{equation}\tag{6.2}
f_*:K_\bullet(X)\longrightarrow K_\bullet(Y)
\end{equation}
est compatible avec les filtrations, et d\'efinit un homomorphisme de groupes gradu\'es
\begin{equation}\tag{6.3}
f_*:\Gr_\bullet(X)\longrightarrow \Gr_\bullet(Y),
\end{equation}
donnant lieu \`a une «~formule de projection~» i.e. qui est $\Gr^\bullet(Y)$-lin\'eaire, formule qui se r\'eduit par passage aux gradu\'es de la formule de projection analogue pour l'homomorphisme (6.2) (qui est $K^\bullet(Y)$-lin\'eaire).

\uline{Dans le point de vue propos\'e dans le S\'eminaire}, \uline{une th\'eorie des intersections maniable}, \uline{sur les sch\'emas noeth\'eriens} $X$ \uline{pas n\'ecessairement r\'eguliers}, \uline{doit consister en l'\'etude combin\'ee des invariants contravariants} $K^\bullet(X),\Gr^\bullet(X)$ \uline{et des invariants covariants} $K_\bullet(X),\Gr_\bullet(X)$, \uline{qui jouent des r\^oles de nature essentiellement diff\'erente}, \uline{et dont les propri\'et\'es se compl\`etent mutuellement}. C'est ainsi que les groupes covariants $K_\bullet(X),\Gr_\bullet(X)$ se calculent mieux pour certaines fibrations non propres, gr\^ace \`a la «~suite exacte d'homotopie~» qu'on peut \'etablir pour eux ([RRR] et Exp.\ IX)~; par exemple on prouve (Exp.\ IX) des isomorphismes canoniques tels que
\[
K_\bullet(X[t])\simeq K_\bullet(X),\qquad \Gr_\bullet(X[t])\simeq \Gr_\bullet(X),
\]
qui tombent en d\'efaut en g\'en\'eral pour $K^\bullet$ lorsque $X$ est un sch\'ema non r\'egulier. En revanche, $K^\bullet(X)$ et $\Gr^\bullet(X)$ se pr\^etent bien aux calculs gr\^ace \`a leur structure d'anneaux.

Lorsqu'on consid\`ere des sch\'emas $X$ propres sur un corps, la projection
\[
f^X:X\longrightarrow(\operatorname{point})
\]
d\'efinit un homomorphisme «~\uline{caract\'eristique d'Euler--Poincar\'e}~»
\begin{equation}\tag{6.4}
\lambda_X=f^X_*:K_\bullet(X)\longrightarrow K_\bullet(\operatorname{point})=\mathbb Z,
\end{equation}
induisant sur le sous-groupe $\Gr_0(X)=\operatorname{Filt}_0K_\bullet(X)$ l'homomorphisme «~degr\'e des $0$-cycles~»
\begin{equation}\tag{6.5}
\deg_X:\Gr_0(X)\longrightarrow \mathbb Z.
\end{equation}
L'homomorphisme de multiplication
\[
\Gr^i(X)\times\Gr_i(X)\longrightarrow \Gr_0(X),
\]
compos\'e avec l'homomorphisme degr\'e (6.5), fournit alors un accouplement
\begin{equation}\tag{6.6}
\Gr^i(X)\times \Gr_i(X)\longrightarrow \mathbb Z,
\end{equation}
qui explique qu'\`a certains \'egards $\Gr^\bullet(X)$ et $\Gr_\bullet(X)$ jouent des r\^oles duals,
